\section{Overall Findings and Limitations}

The great advantage of RL is that an agent is able to learn and to perform actions in a
dynamically changing environment. It thus allows the best possible decisions to be made
during specific predefined points in time. The ability to train and improve the policy
independently through several iteration steps therefore allows an optimal applicability in
the field of algorithmic trading strategies. At the same time, the trading environment is
characterized by more or less external influences such as the size and thus the power of
larger investors, the influence of news and daily developments, decisions of individual
investors, etc. It can further be assumed that the generalizability of an RL model is
difficult to realize and it can therefore be a challenge to derive clear future behavioral
implications.

Applied to my example, I could show that, for example, the application of the A2C
algorithm to the Boeing stock over the last 10 years would lead to an average negative
reward compared to a simple buy-and-hold strategy over the same period. This could be
explained by the fact that especially drastic external influences have led the stock to
behave in an unusual way. Out of 5 trained agents, one in particular showed very poor
results of less than 1, which supports my thesis of the time dependency and volatility of
the model. Nevertheless, for the other two stocks of Airbus and Toyota over a time period
of 10 years, the algorithm showed significantly better results in some cases, but without
outputting consistent behavior once the agent is trained several times. It is further
interesting to note that results in terms of positive rewards do not deviate much when all
technical indicators are used on the one hand and only the 11 indicators that correlate
most with tomorrow's log-return of the close price on the other.

By training the agent on the Toyota stock and then testing the model over a different time
period, I was able to show that the algorithm used on this stock is unlikely to generate
promising results in the future.
